\documentclass[11pt]{article}

\usepackage[letterpaper,margin=1in]{geometry}
\usepackage{amsmath,amssymb,amsthm,mathtools,bm}
\usepackage{booktabs}
\usepackage{enumitem}
\usepackage[colorlinks=true,linkcolor=blue,urlcolor=blue]{hyperref}
\usepackage{microtype}

\numberwithin{equation}{section}

\newtheorem{theorem}{Theorem}[section]
\newtheorem{lemma}[theorem]{Lemma}
\newtheorem{corollary}[theorem]{Corollary}
\newtheorem{remark}[theorem]{Remark}

\DeclareMathOperator{\Tr}{Tr}
\DeclareMathOperator{\diag}{diag}
\newcommand{\R}{\mathbb R}
\newcommand{\E}{\mathbb E}
\newcommand{\Id}{\mathrm{Id}}
\newcommand{\op}{\mathrm{op}}
\newcommand{\loc}{\mathrm{loc}}

\title{A Smooth Genuine-Hessian Counterexample to the\\
Dimension-Free Logarithmic Local-Rate Conjecture}
\author{}
\date{July 14, 2026}

\begin{document}

\maketitle

\begin{abstract}
We construct an explicit sequence of smooth, strongly convex potentials in
equilibrium-whitened coordinates for which the Gaussian natural-gradient
linearized gap has the square-root, rather than logarithmic, condition-number
scale.  In dimension $N_\theta=d+1$, the potentials satisfy
\[
  \E\nabla V_d(Z)=0,\qquad \E\nabla^2V_d(Z)=I,
  \qquad
  \alpha_d I\preceq\nabla^2V_d\preceq\beta_d I,
\]
with post-whitening condition number $\kappa_d=\beta_d/\alpha_d\asymp d$,
but
\[
  \gamma_{\loc}(V_d)\asymp d^{-1/2}\asymp\kappa_d^{-1/2}.
\]
Thus $\gamma_{\loc}^{-1}\asymp\sqrt{\kappa_d}\gg1+\log\kappa_d$.
The same examples satisfy $\tau_H^2\gtrsim\sqrt d\asymp\sqrt{\kappa_d}$,
so they also disprove the proposed dimension-free first-Hermite Hessian
bound.  The Hessian field is pointwise $\nabla^2V_d$ with
$V_d\in C^\infty$; no symmetry-only surrogate, noncommutative Loewner
transfer, or unproved Bellman/packet input is used.
\end{abstract}

\section{Statement of the counterexample}

All constants below are numerical and independent of $d$.  The notation
$a_d\asymp b_d$ means $c b_d\le a_d\le C b_d$ for two such constants.

\begin{theorem}[Smooth square-root counterexample]\label{thm:main}
There are universal constants $c,C>0$ and $d_0\in\mathbb N$ such that, for
every integer $d\ge d_0$, there is a $C^\infty$, strongly convex potential
$V_d:\R^{d+1}\to\R$ with Hessian $W_d=\nabla^2V_d$ satisfying, for
$Z\sim\mathcal N(0,I_{d+1})$,
\begin{equation}\label{eq:equilibrium}
  \E\nabla V_d(Z)=0,
  \qquad
  \E W_d(Z)=I_{d+1}.
\end{equation}
If $\alpha_d$ and $\beta_d$ are the optimal scalar lower and upper curvature
constants, then
\begin{equation}\label{eq:kappa-scale}
  \alpha_d=d^{-1/2},
  \qquad
  c\sqrt d\le\beta_d\le C\sqrt d,
  \qquad
  c d\le\kappa_d:=\frac{\beta_d}{\alpha_d}\le C d.
\end{equation}
The positive linearized generator at $(0,I)$ obeys
\begin{equation}\label{eq:gap-scale}
  \frac1{\sqrt d}
  \le\gamma_{\loc}(V_d)
  \le\frac C{\sqrt d}.
\end{equation}
Consequently
\begin{equation}\label{eq:counter-conclusion}
  \gamma_{\loc}(V_d)\asymp\kappa_d^{-1/2},
  \qquad
  \gamma_{\loc}(V_d)^{-1}\ge c\sqrt{\kappa_d},
  \qquad
  (1+\log\kappa_d)\gamma_{\loc}(V_d)\longrightarrow0.
\end{equation}
Moreover, the first-Hermite Hessian coupling satisfies
\begin{equation}\label{eq:tau-lower-main}
  \tau_H(V_d)^2\ge c\sqrt d\ge c'\sqrt{\kappa_d}.
\end{equation}
\end{theorem}

The construction is given explicitly below.  Its only dimension-dependent
parameter $\delta_d$ is the smallest zero in $[0,1]$ of the continuous scalar
function in \eqref{eq:isotropizing-equation}; Lemma~\ref{lem:delta} proves
that this zero exists and has the stated scale.  This is a canonical
one-dimensional definition, analogous to specifying a Gaussian quantile.

\section{The radial-softplus Hessian}

Write $z=(t,y)\in\R\times\R^d$ and define
\begin{equation}\label{eq:profiles}
  r(y):=\sqrt{1+\lVert y\rVert^2},
  \qquad
  F(s):=\log(1+e^s),
  \qquad
  p(s):=F'(s),
  \qquad
  q(s):=F''(s)=p(s)(1-p(s)).
\end{equation}
Thus $0<p<1$ and $0<q\le1/4$.  For $0\le\delta\le1$, set
\begin{equation}\label{eq:Udelta}
  s_\delta(t,y):=r(y)-\sqrt d-\delta t,
  \qquad
  U_\delta(t,y):=F(s_\delta(t,y)),
  \qquad
  A_\delta:=\nabla^2U_\delta.
\end{equation}
Let $S=\lVert y\rVert^2$ and
\begin{equation}\label{eq:radial-derivatives}
  v_\delta:=\nabla s_\delta=\begin{pmatrix}-\delta\\y/r\end{pmatrix},
  \qquad
  D(y):=\nabla_y^2r=\frac{I_d}{r}-\frac{yy^\top}{r^3}.
\end{equation}
Direct differentiation gives the exact Hessian
\begin{equation}\label{eq:HessU}
  A_\delta
  =q(s_\delta)v_\delta v_\delta^\top
  +p(s_\delta)
   \begin{pmatrix}0&0\\0&D(y)\end{pmatrix}.
\end{equation}
In longitudinal/transverse blocks this is
\begin{equation}\label{eq:Hess-blocks}
  (A_\delta)_{tt}=\delta^2q,
  \qquad
  (A_\delta)_{yt}=-\delta q\frac y r,
  \qquad
  (A_\delta)_{yy}=q\frac{yy^\top}{r^2}+pD(y).
\end{equation}
Both summands in \eqref{eq:HessU} are positive semidefinite.  The eigenvalues
of $D$ are $r^{-1}$ in tangential directions and $r^{-3}$ in the radial
direction.  Hence $0\preceq D\preceq I_d$, while
$\lVert v_\delta\rVert^2\le2$.  It follows that
\begin{equation}\label{eq:A-bound}
  0\preceq A_\delta\preceq\frac32 I_{d+1}
  \qquad(0\le\delta\le1).
\end{equation}

Let $T\sim\mathcal N(0,1)$ and $Y\sim\mathcal N(0,I_d)$ be independent.  In
the next definitions, $s_\delta=s_\delta(T,Y)$ and $r=r(Y)$:
\begin{align}
  Q_d(\delta)
  &:=\E q(s_\delta),\label{eq:Qdef}\\
  M_d(\delta)
  &:=\frac1d\E\left[
       q(s_\delta)\frac{S}{r^2}
       +p(s_\delta)\left(\frac d r-\frac S{r^3}\right)
     \right].\label{eq:Mdef}
\end{align}
Rotation invariance in $y$ and oddness of the mixed block imply
\begin{equation}\label{eq:mean-A}
  \E A_\delta
  =\diag\bigl(\delta^2Q_d(\delta),M_d(\delta)I_d\bigr).
\end{equation}

\begin{lemma}[Isotropizing parameter]\label{lem:delta}
Uniformly over $0\le\delta\le1$ and all sufficiently large $d$,
\begin{equation}\label{eq:QM-bounds}
  c\le Q_d(\delta)\le\frac14,
  \qquad
  \frac c{\sqrt d}\le M_d(\delta)\le\frac C{\sqrt d}.
\end{equation}
Consequently the continuous equation
\begin{equation}\label{eq:isotropizing-equation}
  \delta^2Q_d(\delta)=M_d(\delta)
\end{equation}
has a zero $\delta_d\in(0,1)$.  If $\delta_d$ denotes its smallest zero and
\begin{equation}\label{eq:mdef}
  m_d:=\delta_d^2Q_d(\delta_d)=M_d(\delta_d),
\end{equation}
then
\begin{equation}\label{eq:delta-m-scale}
  c d^{-1/4}\le\delta_d\le C d^{-1/4},
  \qquad
  c d^{-1/2}\le m_d\le C d^{-1/2},
  \qquad
  \E A_{\delta_d}=m_dI_{d+1}.
\end{equation}
\end{lemma}

\begin{proof}
Let
\[
  \mathcal E_d:=\{|S-d|\le2\sqrt d,\ |T|\le1\}.
\]
Since $\operatorname{Var}(S)=2d$, Chebyshev's inequality and independence
give $\mathbb P(\mathcal E_d)\ge c_0>0$.  On $\mathcal E_d$,
\[
  (\sqrt d-1)^2\le1+S\le(\sqrt d+1)^2,
\]
so $|r-\sqrt d|\le1$ and $|s_\delta|\le2$ for every
$0\le\delta\le1$.  Therefore $q(s_\delta)\ge\min_{|s|\le2}q(s)>0$ and
$p(s_\delta)\ge p(-2)>0$ on this event.  This proves the lower bound on
$Q_d$; the upper bound is $q\le1/4$.

The identity
\begin{equation}\label{eq:trace-D}
  \frac d r-\frac S{r^3}=\frac{d-1}{r}+\frac1{r^3}
\end{equation}
shows that the second term in \eqref{eq:Mdef} is at least $c\sqrt d$ on
$\mathcal E_d$.  This proves $M_d(\delta)\ge c/\sqrt d$.

For the upper bound, the first integrand in \eqref{eq:Mdef} is at most
$1/4$.  Also, for $d>2$,
\begin{equation}\label{eq:inverse-radius}
  \E\frac1r
  \le\E S^{-1/2}
  \le(\E S^{-1})^{1/2}
  =\frac1{\sqrt{d-2}},
\end{equation}
where the last equality follows by integrating the $\chi_d^2$ density.
Combining \eqref{eq:trace-D} and \eqref{eq:inverse-radius} gives
$M_d(\delta)\le C/\sqrt d$.

The map $\delta\mapsto\delta^2Q_d(\delta)-M_d(\delta)$ is continuous by
dominated convergence.  It is negative at $\delta=0$ and, by
\eqref{eq:QM-bounds}, positive at $\delta=1$ for all sufficiently large $d$.
Thus a zero exists.  At every zero, \eqref{eq:QM-bounds} gives
$\delta_d^2\asymp d^{-1/2}$ and $m_d\asymp d^{-1/2}$.  The last assertion
then follows from \eqref{eq:mean-A}.
\end{proof}

\section{Whitening and the post-whitening condition number}

Fix $\delta_d$ from Lemma~\ref{lem:delta}, abbreviate
$U_d:=U_{\delta_d}$ and $A_d:=A_{\delta_d}$, and set
\begin{equation}\label{eq:alpha-c}
  \alpha_d:=d^{-1/2},
  \qquad
  c_d:=\frac{1-\alpha_d}{m_d}.
\end{equation}
Define
\begin{equation}\label{eq:Vdef}
  V_d(z):=\frac{\alpha_d}{2}\lVert z\rVert^2+c_dU_d(z)+b_d^\top z,
  \qquad
  b_d:=-c_d\E\nabla U_d(Z)
      =c_d\delta_d\E p(s_{\delta_d}(Z))\,e_t.
\end{equation}
The vector $b_d$ is finite because $\nabla U_d$ is bounded; explicitly, it
is a multiple of the $t$ direction.  It changes no Hessian and makes
$\E\nabla V_d=0$.  The Hessian is
\begin{equation}\label{eq:Wdef}
  W_d(z)=\nabla^2V_d(z)=\alpha_dI+c_dA_d(z).
\end{equation}
By \eqref{eq:delta-m-scale} and \eqref{eq:alpha-c},
\begin{equation}\label{eq:whitening}
  \E W_d=\alpha_dI+c_dm_dI=I.
\end{equation}
Thus the construction is already in equilibrium-whitened coordinates; every
condition number below is $\kappa_{\rm white}$.

From \eqref{eq:A-bound},
\begin{equation}\label{eq:W-upper}
  \alpha_dI\preceq W_d
  \preceq\left(\alpha_d+\frac32c_d\right)I.
\end{equation}
For fixed $y$, $A_d(t,y)\to0$ as $t\to+\infty$, so the optimal lower
curvature constant is exactly $\alpha_d$.  At $y=0$ and
$t=(1-\sqrt d)/\delta_d$, one has $s_{\delta_d}=0$ and the transverse block
of $A_d$ is $I_d/2$.  Hence the optimal upper constant $\beta_d$ satisfies
\begin{equation}\label{eq:beta-bounds}
  \alpha_d+\frac12c_d\le\beta_d
  \le\alpha_d+\frac32c_d.
\end{equation}
Since $m_d\asymp d^{-1/2}$, this proves
$\beta_d\asymp\sqrt d$ and $\kappa_d=\beta_d/\alpha_d\asymp d$.

The function $V_d$ is $C^\infty$ and $\alpha_d$-strongly convex.  In
particular $e^{-V_d}$ is integrable, and $W_d$ is a genuine pointwise Hessian
whose complete mixed-derivative hierarchy holds classically (and hence also
distributionally).

\section{The exact local Rayleigh identity}

For completeness, we record the coordinate convention used for the local
gap.  If $X\in\operatorname{Sym}(d+1)$ is the physical covariance
perturbation, its packed coordinate is $X/\sqrt2$ and its Fisher--Rao norm is
\begin{equation}\label{eq:FRnorm}
  \lVert(u,X)\rVert_\star^2
  :=\lVert u\rVert^2+\frac12\lVert X\rVert_F^2.
\end{equation}

\begin{lemma}[Compatible Bochner--Rayleigh identity]\label{lem:rayleigh}
For a smooth Hessian field $W=\nabla^2V$ with $\E W=I$, the quadratic form of
the packed generator in the question satisfies
\begin{equation}\label{eq:rayleigh-identity}
  4\mathcal Q_W(u,X)
  =\E\big[(2u+XZ)^\top W(Z)(2u+XZ)\big]+\lVert X\rVert_F^2,
\end{equation}
where
$\mathcal Q_W(u,X)=\langle(u,X/\sqrt2),L_\star(u,X/\sqrt2)\rangle$.
Consequently, if $W\succeq\alpha I$ with $0<\alpha\le1$, then
\begin{equation}\label{eq:alpha-floor}
  \mathcal Q_W(u,X)\ge
  \alpha\left(\lVert u\rVert^2+\frac12\lVert X\rVert_F^2\right),
  \qquad
  \gamma_{\loc}(W)\ge\alpha.
\end{equation}
\end{lemma}

\begin{proof}
Expanding the displayed block operator $L_\star$ in the packed coordinate
$(u,X/\sqrt2)$ gives
\[
  \mathcal Q_W(u,X)
  =\lVert u\rVert^2+u^\top T[X]
   +\frac12\lVert X\rVert_F^2+\frac14\langle X,H[X]\rangle_F.
\]
Gaussian integration by parts, together with
$\partial_kW_{ij}=\partial_iW_{kj}$ and the corresponding fourth-order
mixed-derivative symmetry, gives the compatible Bochner identities
\begin{align*}
  \E[u^\top W XZ]&=u^\top T[X],\\
  \E[(XZ)^\top W(XZ)]
    &=\lVert X\rVert_F^2+\langle X,H[X]\rangle_F.
\end{align*}
Substitution proves \eqref{eq:rayleigh-identity}.  All integrations by parts
are classical for $V_d$; equivalently they follow first for compactly
supported cutoffs and then by dominated convergence from the bounded Hessian.

If $W\succeq\alpha I$, then $\E Z=0$ and
$\E\lVert XZ\rVert^2=\lVert X\rVert_F^2$ give
\[
  4\mathcal Q_W(u,X)
  \ge4\alpha\lVert u\rVert^2+(1+\alpha)\lVert X\rVert_F^2
  \ge4\alpha\lVert(u,X)\rVert_\star^2,
\]
where the last step uses $\alpha\le1$.
\end{proof}

\section{Explicit inverse-gap certificate}

Take
\begin{equation}\label{eq:witness}
  u=e_t,
  \qquad
  X=\diag(0,\lambda_d I_d),
  \qquad
  \lambda_d:=\frac{2\delta_d}{\sqrt d}.
\end{equation}
Then
\begin{equation}\label{eq:b-and-X}
  2u+XZ=(2,\lambda_dY),
  \qquad
  \lVert X\rVert_F^2=d\lambda_d^2=4\delta_d^2,
  \qquad
  \lVert(u,X)\rVert_\star^2=1+2\delta_d^2\ge1.
\end{equation}
Put
\begin{equation}\label{eq:gdef}
  g(Y):=\frac{\lVert Y\rVert^2}{r(Y)}=r(Y)-\frac1{r(Y)}.
\end{equation}
Using \eqref{eq:HessU}, the Hessian energy is exactly
\begin{equation}\label{eq:exact-A-energy}
  (2,\lambda_dY)^\top A_d(2,\lambda_dY)
  =q(s_{\delta_d})\bigl(-2\delta_d+\lambda_dg(Y)\bigr)^2
   +p(s_{\delta_d})\lambda_d^2\frac{S}{r^3}.
\end{equation}

Two elementary radial bounds suffice.  First,
\begin{equation}\label{eq:r-centered}
  r-\sqrt d=\frac{S+1-d}{r+\sqrt d}
  \quad\Longrightarrow\quad
  \E(r-\sqrt d)^2\le\frac{\E(S+1-d)^2}{d}le3.
\end{equation}
Since $g=r-r^{-1}$ and $r\ge1$,
\begin{equation}\label{eq:g-centered}
  \E(g-\sqrt d)^2\le8.
\end{equation}
Second, \eqref{eq:inverse-radius} gives
\begin{equation}\label{eq:radial-Hess-mean}
  \E\frac{S}{r^3}\le\E\frac1r\le\frac1{\sqrt{d-2}}.
\end{equation}
Because $q\le1/4$ and $\lambda_d=2\delta_d/\sqrt d$,
\eqref{eq:exact-A-energy}--\eqref{eq:radial-Hess-mean} imply
\begin{equation}\label{eq:A-energy-bound}
  \E[(2,\lambda_dY)^\top A_d(2,\lambda_dY)]
  \le\frac{8\delta_d^2}{d}
      +\frac{4\delta_d^2}{d\sqrt{d-2}}
  \le\frac{C\delta_d^2}{d}.
\end{equation}
At the isotropizing root, $m_d=\delta_d^2Q_d(\delta_d)$ and
$Q_d(\delta_d)\ge c$.  Therefore
\begin{equation}\label{eq:scaled-A-energy}
  c_d\E[(2,\lambda_dY)^\top A_d(2,\lambda_dY)]
  \le\frac Cd.
\end{equation}
The remaining two terms in \eqref{eq:rayleigh-identity} satisfy
\begin{equation}\label{eq:floor-and-X-cost}
  \alpha_d\E\lVert(2,\lambda_dY)\rVert^2
  =\alpha_d(4+4\delta_d^2)\le\frac C{\sqrt d},
  \qquad
  \lVert X\rVert_F^2=4\delta_d^2\le\frac C{\sqrt d}.
\end{equation}
Combining \eqref{eq:rayleigh-identity}, \eqref{eq:scaled-A-energy}, and
\eqref{eq:floor-and-X-cost}, then dividing by
\eqref{eq:b-and-X}, gives
\begin{equation}\label{eq:gap-upper}
  \gamma_{\loc}(V_d)le\frac C{\sqrt d}.
\end{equation}
The matching lower bound $\gamma_{\loc}\ge\alpha_d=d^{-1/2}$ follows from
\eqref{eq:alpha-floor}.  Together with $\kappa_d\asymp d$, this proves
\eqref{eq:gap-scale}--\eqref{eq:counter-conclusion}.

Geometrically, the witness is not accidental.  The rank-one vector in
\eqref{eq:HessU} is $(-\delta_d,Y/r)$, and its inner product with
$(2,\lambda_dY)$ is
$2\delta_d(g/\sqrt d-1)$.  Gaussian radial concentration makes this
$O_{L^2}(\delta_d/\sqrt d)$, while the radial eigenvalue of $D$ is $r^{-3}$.
The affine perturbation therefore follows an approximate null direction of the
\emph{genuine} Hessian.

\section{First-Hermite certificate}

Let
\begin{equation}\label{eq:Btest}
  B:=\diag(0,I_d/\sqrt d),
  \qquad
  \lVert B\rVert_F=1.
\end{equation}
For the operator $T$ in the question, smooth Hessian compatibility and two
Gaussian integrations by parts give
\begin{align}
  e_t^\top T[B]
  &=\E[t\Tr(BW_d)]
    =\E[(BZ)^\top W_de_t]\notag\\
  &=-c_d\frac{\delta_d}{\sqrt d}
    \E\left[q(s_{\delta_d})\frac{S}{r}\right].
  \label{eq:T-certificate}
\end{align}
On the event $\mathcal E_d$ used in Lemma~\ref{lem:delta},
$q(s_{\delta_d})\ge c$ and $S/(\sqrt d\,r)\ge c$ for all sufficiently
large $d$.  Hence
\begin{equation}\label{eq:qzeta-lower}
  \frac1{\sqrt d}\E\left[q(s_{\delta_d})\frac S r\right]\ge c.
\end{equation}
Since $c_d\asymp\sqrt d$ and $\delta_d\asymp d^{-1/4}$,
\eqref{eq:T-certificate}--\eqref{eq:qzeta-lower} yield
\begin{equation}\label{eq:tau-lower}
  \tau_H^2\ge|e_t^\top T[B]|^2\ge c\sqrt d.
\end{equation}
Because $\kappa_d\asymp d$,
\[
  D_{\kappa_d}
  =\left(4+\frac4{\sqrt\pi}\right)(1+\log\kappa_d)
  =O(\log d)
  =o(\tau_H^2).
\]
This proves \eqref{eq:tau-lower-main} and directly refutes the intended
dimension-free first-Hermite estimate.

\begin{corollary}[The affine weighted Korn frontier also fails]
Put $\widetilde A_d=W_d-\alpha_dI=c_dA_d$.  Then
$\E\widetilde A_d=(1-\alpha_d)I$ and
$0\preceq\widetilde A_d\preceq L_dI$ with $L_d\asymp\sqrt d$.  For the
witness \eqref{eq:witness},
\[
  \E[(2u+XZ)^\top\widetilde A_d(2u+XZ)]\le\frac Cd,
  \qquad
  (1-\alpha_d)\lVert X\rVert_F^2\le\frac C{\sqrt d}.
\]
Thus multiplying their sum by
$1+\log(eL_d/(1-\alpha_d))=O(\log d)$ still gives $o(1)$, whereas
$(1-\alpha_d)\lVert u\rVert^2\to1$.  Hence the proposed affine weighted Korn
inequality cannot hold under the stated assumptions, even for smooth genuine
Hessians.
\end{corollary}

\section{Guardrail audit}

\begin{enumerate}[leftmargin=2em,label=\arabic*.]
  \item The construction never replaces Hessian compatibility by symmetry of
  the first Hermite tensor.  Equation \eqref{eq:HessU} is the pointwise
  Hessian of the displayed $C^\infty$ potential.
  \item No commutation-free Loewner transfer is used.  The inverse-gap
  certificate is the scalar Rayleigh computation
  \eqref{eq:exact-A-energy}--\eqref{eq:gap-upper}.
  \item Spectral bounds and isotropic mean are not treated as sufficient
  abstract data.  The small energy comes from the explicit compatible radial
  and mixed Hessian blocks in \eqref{eq:Hess-blocks}.
  \item Whitening is exact in \eqref{eq:whitening}.  The constants
  $\alpha_d,\beta_d,\kappa_d$ are post-whitening constants; no original-to-
  whitened conversion is invoked.
  \item The potential is $C^\infty$, its Hessian is bounded, and all Gaussian
  integrations by parts and mixed-derivative identities are classical.  They
  also remain valid in the weak/distributional formulation.
  \item The construction is genuinely high-dimensional and gives
  $\gamma_{\loc}^{-1}\asymp\sqrt\kappa$, so it does not conflict with the
  sharp one-dimensional logarithmic examples.
  \item The result concerns only the linearized local gap.  It makes no claim
  about a nonlinear attraction radius.
\end{enumerate}

\paragraph{Conclusion.}
The family \eqref{eq:Vdef} is an admissible option-(B) counterexample.  The
dimension-free logarithmic local-rate conjecture is false under the assumptions
in the question, and the already-known square-root fallback has the correct
order on this smooth family.

\end{document}
